Promotions can build customer engagement while lowering the reference prices against which future offers are judged. When should a seller spread promotions evenly, and when should it concentrate them into campaigns? We develop a common framework for these decisions with fixed discount depth and finite customer memory. Demand may depend nonlinearly on customer age, current price, and the number of remembered promotions, allowing both positive and negative historical effects. Our central insight is that the shape of window contribution functions derived from demand gives a common criterion for optimal arrangement and tractable frequency choice. Concavity makes evenly spread arrangements optimal at a given cycle length and promotion count and reduces frequency choice to a finite search over short cycles. Convexity makes a single promotion block optimal under the same constraints and a constant price optimal when frequency is unrestricted. \textcolor{blue}{Outside these regimes, arrangement at prescribed cycle length and promotion count is NP-hard even for affine demand;} simple policies and memory truncation have revenue-loss guarantees. In a synthetic access-service example, spreading the same promotion budget raises revenue by 2.32\% relative to clustering and substantially smooths demand.
